\documentclass[11pt]{article}
\usepackage[margin=1in]{geometry}
\usepackage{amsmath,amssymb,amsthm,booktabs,array}
\usepackage[T1]{fontenc}
\usepackage{lmodern}
\usepackage{microtype}
\usepackage[hidelinks]{hyperref}
\setlength{\emergencystretch}{2em}
\newtheorem{theorem}{Theorem}
\newtheorem{lemma}{Lemma}
\newcommand{\R}{\mathbb R}
\newcommand{\Sph}{\mathbb S^2}
\newcommand{\ip}[2]{\langle #1,#2\rangle}
\newcommand{\HH}{\mathcal H^2}
\newcommand{\code}[1]{\texttt{\detokenize{#1}}}
\title{Minimum spherical 5-designs on the two-sphere\\\large Conjecture 00000001030: minimum, equality, and the exception}
\author{C0ldSmi1e}
\date{October 12, 2026}
\begin{document}
\maketitle
\begin{abstract}
Every nonempty finite spherical 5-design on the Euclidean unit two-sphere has at least twelve points. The normalized vertices of a regular icosahedron attain twelve, and every twelve-point design is orthogonally congruent to that set. Thus the ordinary minimum is twelve. If the original parenthetical is instead read as excluding icosahedral designs, twelve is impossible: every remaining design has at least thirteen points. We make no claim that thirteen is attained. All of these statements are proved in Lean against normalized two-dimensional Hausdorff surface measure, including the measure and polynomial-exactness bridges.
\end{abstract}

\section{Original statement and precise scope}
The supplied English conjecture is:
\begin{quote}
The minimal number of points of a spherical 5-design on S\textsuperscript{2} is 12 (with the icosahedron exceptional) (minimal 5-designs, Bannai's count).
\end{quote}
The accompanying Chinese statement also places the icosahedron in a parenthetical exception. The exact bilingual source is included as \code{conjecture.md}. To avoid silently choosing the meaning of that parenthetical, we establish both conclusions described in the abstract.

Let $E=\R^3$ with its Euclidean inner product, and let
$\Sph=\{x\in E:\|x\|_2=1\}$. Write $\sigma$ for normalized Hausdorff surface measure on $\Sph$. A spherical 5-design is a nonempty finite \emph{set} $X\subset\Sph$ such that, for every real polynomial $p$ in three variables of total degree at most five,
\begin{equation}\label{eq:design}
 \frac1{|X|}\sum_{x\in X}p(x)=\int_{\Sph}p(u)\,d\sigma(u).
\end{equation}
Points are distinct and all weights are equal. Set
\[
 \varphi=\frac{1+\sqrt5}{2},\qquad r=\sqrt{\varphi+2},\qquad
 a=\frac1r,\qquad b=\frac\varphi r,
\]
and define the twelve-point coordinate set
\begin{equation}\label{eq:ico}
 I=\{(0,\pm a,\pm b),(\pm a,\pm b,0),(\pm b,0,\pm a)\},
\end{equation}
where the signs are independent in each family. This is the standard normalized icosahedral vertex set. In the formal statement, ``icosahedral'' means exactly that an orthogonal linear isometry carries $X$ onto $I$.

\begin{theorem}\label{thm:main}
Every spherical 5-design $X$ satisfies $|X|\ge12$. The set $I$ is such a design and has twelve points. For any spherical 5-design $X$,
\[
 |X|=12\quad\Longleftrightarrow\quad \exists U\in O(3),\ U(X)=I.
\]
Consequently, no nonicosahedral spherical 5-design has twelve points; each has at least thirteen.
\end{theorem}
The last conclusion neither asserts existence of a nonicosahedral design nor identifies that class's exact minimum. It refutes an \emph{attained minimum of twelve} under the exclusion reading.

\section{Surface measure and moments}
The Lean space is \code{EuclideanSpace} $\R$ \code{(Fin 3)}, so its norm is the Euclidean norm. We take the intrinsic Hausdorff measure $\HH$ on the unit-sphere subtype. The inclusion into $E$ is an isometry, and its pushforward is ambient $\HH$ restricted to the sphere. This bridge is proved rather than assumed.

Here are explicit proofs that its total mass is positive and finite. On the planar unit disk, use the two maps
\[
 s_\varepsilon(u)=\frac{(2u_1,2u_2,\varepsilon(1-\|u\|^2))}{1+\|u\|^2},
 \qquad \varepsilon\in\{1,-1\}.
\]
They take values on $\Sph$; their images cover its two closed hemispheres, including the equator. The identity
\[
 \|s_\varepsilon(u)-s_\varepsilon(v)\|^2
 =\frac{4\|u-v\|^2}{(1+\|u\|^2)(1+\|v\|^2)}
\]
shows that each is 2-Lipschitz. The planar disk has finite two-dimensional Hausdorff measure, giving finiteness of the sphere's measure. Conversely, orthogonal projection of the sphere onto its first two coordinates is 1-Lipschitz and has image the closed unit disk, whose measure is positive. The Hausdorff image inequality gives positivity of the sphere's measure.

Thus normalization by this finite positive mass defines a probability measure $\sigma$. Orthogonal maps preserve Hausdorff measure and the sphere, hence preserve $\sigma$. Restricted polynomials are continuous on the compact sphere and therefore integrable. These facts are all proved in \code{SurfaceMeasure.lean}; no unproved numerical surface area is needed.

Coordinate sign changes show that any monomial having an odd exponent has integral zero. Coordinate permutations and $u_1^2+u_2^2+u_3^2=1$ give
\[
 \int u_i^2\,d\sigma=\frac13.
\]
For fourth moments, put $A=\int u_1^4\,d\sigma$ and
$B=\int u_1^2u_2^2\,d\sigma$. Permutation invariance makes these independent of the selected distinct coordinates. The orthogonal coordinate change
\[
 (u_1,u_2,u_3)\longmapsto
 ((u_1+u_2)/\sqrt2,(u_1-u_2)/\sqrt2,u_3)
\]
gives $A=(2A+6B)/4$, so $A=3B$. Multiplying the norm identity by $u_1^2$ and integrating gives $A+2B=1/3$. Therefore
\begin{equation}\label{eq:moments}
 \int u_i^4\,d\sigma=\frac15,\qquad
 \int u_i^2u_j^2\,d\sigma=\frac1{15}\quad(i\ne j).
\end{equation}
The displayed coordinate change is an orthogonal reflection, not a proper rotation. Together with the constant moment one and the odd-exponent rule, these formulas exhaust every monomial of total degree at most five. Finite expansion over the support of a multivariate polynomial proves the bridge from these monomial identities to \eqref{eq:design}.

For any fixed unit vector $x$, an orthogonal reflection carries $x$ to the first coordinate vector. Consequently the dot-product moments are
\begin{equation}\label{eq:dotmoments}
 \int\ip{x}{u}^{\,k}\,d\sigma(u)=
 \begin{cases}1&k=0,\\1/3&k=2,\\1/5&k=4,\\0&k=1,3,5.\end{cases}
\end{equation}
The Lean polynomial \code{linearForm x} represents $u\mapsto\ip{x}{u}$, with a proved degree bound. Thus design exactness applies to each of its first five powers without adding a separate moment hypothesis to the final theorem.

\section{Lower bound and equality constraints}
Consider the degree-five polynomial
\[
 F(t)=(1+t)(5t^2-1)^2
     =25t^5+25t^4-10t^3-10t^2+t+1.
\]
It is nonnegative on $[-1,1]$, and $F(1)=32$. By \eqref{eq:dotmoments}, its sphere average at any unit pole is
\[
 25\cdot\frac15-10\cdot\frac13+1=\frac83.
\]
For a design of size $N$ and a point $x\in X$, exactness therefore gives
\begin{equation}\label{eq:lp}
 \frac83N=\sum_{y\in X}F(\ip{x}{y})\ \ge\ F(1)=32.
\end{equation}
The inner products lie in $[-1,1]$ by Cauchy--Schwarz, so every summand used in the inequality is nonnegative. This proves $N\ge12$.

If $N=12$, the sum equals its diagonal term. Hence every off-diagonal term is zero, and distinct $x,y\in X$ satisfy
\begin{equation}\label{eq:roots}
 \ip{x}{y}=-1\quad\hbox{or}\quad \ip{x}{y}^2=\frac15.
\end{equation}
These constraints also imply antipodality. If $-x$ were absent, all eleven off-diagonal squared inner products would equal $1/5$, giving
\[
 \sum_{y\in X}\ip{x}{y}^2=1+\frac{11}{5}=\frac{16}{5}.
\]
But degree-two exactness gives $12/3=4$, a contradiction. The equality case $\ip{x}{y}=-1$ for unit vectors implies $y=-x$, since $\|x+y\|^2=0$.

\section{A twelve-point design}
The identities $\varphi^2=\varphi+1$ yield
\begin{equation}\label{eq:ab}
 a^2+b^2=1,\qquad a^2b^2=\frac15,\qquad
 a^4+b^4=\frac35,\qquad b^2-a^2=ab.
\end{equation}
Both $a$ and $b$ are positive. The three families in \eqref{eq:ico} have different zero coordinates and contain four distinct points each, so $|I|=12$. Every point has unit norm.

Independent coordinate sign changes preserve $I$, so odd-exponent monomial sums vanish. Its cyclic coordinate symmetry and \eqref{eq:ab} give, for every coordinate $i$ and distinct $i,j$,
\[
 \frac1{12}\sum_{v\in I}v_i^2=\frac{4a^2+4b^2}{12}=\frac13,
 \qquad
 \frac1{12}\sum_{v\in I}v_i^4=\frac{4a^4+4b^4}{12}=\frac15,
\]
\[
 \frac1{12}\sum_{v\in I}v_i^2v_j^2=\frac{4a^2b^2}{12}=\frac1{15}.
\]
The constant average is one. These exhaust the even monomials of degree at most five, and match the actual surface moments. The finite polynomial-support bridge now proves that $I$ is a spherical 5-design.

\section{Rigidity of the equality case}
Let $c=ab>0$, so $c^2=1/5$ and $c<1$. By \eqref{eq:roots}, distinct points of a twelve-point design have inner product in $\{-1,c,-c\}$, and the set is antipodal. Choose $x\in X$ and a point outside $\{x,-x\}$, which exists because $|X|=12$. Replacing that point by its antipode if necessary gives $y\in X$ with $\ip{x}{y}=c$.

Put $p=(0,a,b)$ and $q=(0,-a,b)$, for which $\ip{p}{q}=c$. Two orthogonal reflections align $(x,y)$ with $(p,q)$. First reflect in $(x-p)^\perp$ to send $x$ to $p$; then reflect in $(R_1y-q)^\perp$ to send $R_1y$ to $q$. The second reflection fixes $p$ because $\ip{p}{R_1y}=\ip{p}{q}$. The construction also handles a zero difference, in which case the corresponding reflection is the identity.

Write $U$ for their composition. Suppose $z\in X$ has $\ip{x}{z}=c$. If $z=y$, its image is $q\in I$. Otherwise $\ip{y}{z}$ cannot be $-1$: that would give $z=-y$ and hence $\ip{x}{z}=-c$. Thus, writing $w=Uz$,
\[
 \ip{p}{w}=c,\qquad \ip{q}{w}=c\ \hbox{or}\ -c,\qquad \|w\|=1.
\]
Solving the two linear equations and then the norm equation gives exactly
\[
 w=(\pm b,0,a)\quad\hbox{or}\quad w=(\pm a,b,0),
\]
all belonging to $I$. The pole maps to $p$, its antipode maps to $-p$, and every remaining point with pole inner product $-c$ is the negative of one of the preceding upper points. The set $I$ is antipodal, so $U(X)\subseteq I$. An isometry is injective and both sets have twelve points; consequently $U(X)=I$. This proves rigidity and completes Theorem~\ref{thm:main}.

\section{Formal statements and verification}
The project pins Lean 4.19.0 and Mathlib v4.19.0 at commit
\begin{center}\code{c44e0c8ee63ca166450922a373c7409c5d26b00b}.\end{center}
The manifest also pins the transitive dependencies. The main entry point is \code{Conjecture1030.lean}, which imports the final theorem module. The proof is split as follows:
\begin{center}\small
\begin{tabular}{@{}p{0.34\textwidth}p{0.60\textwidth}@{}}
\toprule
Module & Mathematical responsibility\\\midrule
\code{Definitions} & Euclidean sphere, genuine measure and polynomial design definition\\
\code{SurfaceMeasure} & Measure identification, positive finite mass, normalization and invariance\\
\code{SphereMoments} & Actual integrals and the full polynomial-support bridge\\
\code{LinearProgramming} & Nonnegative polynomial, bound, equality roots and antipodes\\
\code{DesignExactness} & Derivation of every finite-sum premise from the actual design definition\\
\code{Icosahedron} & Twelve vertices and all degree-five discrete moments\\
\code{IcosahedronRigidity} & Orthogonal pair alignment and coordinate exhaustion\\
\code{Theorem} & Witness, minimum, rigidity and the exclusion consequence\\\bottomrule
\end{tabular}
\end{center}
The principal declarations, all in namespace \code{Conjecture1030}, are:
\begin{itemize}
\item \code{ico_is_spherical_five_design};
\item \code{spherical_five_design_card_ge_twelve};
\item \code{spherical_five_design_minimum_twelve};
\item \code{spherical_five_design_twelve_rigidity};
\item \code{spherical_five_design_card_twelve_iff};
\item \code{nonicosahedral_spherical_five_design_card_ge_thirteen};
\item \code{no_twelve_point_nonicosahedral_spherical_five_design}.
\end{itemize}
Their statements contain no unproved surface-moment, finite-sum, or geometric-classification premises. Internal lemmas parameterized by invariant probability measures are instantiated with the proved normalized Hausdorff measure before the final capstones.

The author build and strict source replays have succeeded. Axiom queries for the measure bridges and final theorems report only Lean's standard \code{propext}, \code{Classical.choice}, and \code{Quot.sound}. The final independent engineering evidence is recorded separately in the accompanying verification package; it binds results to source hashes and distinguishes a fresh owned-source build from reused generic dependency artifacts. The implementation uses no admitted proofs, \code{native_decide}, or custom axioms.

An independent Python checker uses exact rational arithmetic in $\mathbb Q(\sqrt5)$ to check all 56 monomials of degree at most five, all 66 pairs, and the coordinate model's convex hull. It finds twenty supporting triangular facets and thirty equal-length edges, with the expected incidences. Negative controls remove a vertex, introduce a duplicate, perturb a coordinate by a nonzero rational amount, and alter a moment target; all are rejected. A global reflection and a rational orthogonal coordinate change are accepted. These finite computations supplement the proof; they do not establish the universal bound or replace the formal integral and rigidity arguments. The formal predicate \code{IsIcosahedral} is coordinate congruence, not an additional convex-hull-facet predicate.

From the submission's \code{lean} directory, the standard build and principal axiom queries are:
\begin{verbatim}
lake build
lake env lean -DwarningAsError=true Audit.lean
\end{verbatim}
Run the packaged auxiliary checker with Python 3 according to its README. The included verification instructions give the stricter source replay, dependency-closure audit, and exact output-comparison commands. The complete local verification record is distinct from any later maintainer decision. No claim of a Lean bootstrap rebuild or of rebuilding every Mathlib source is made.
\end{document}
